\section{Compact Hall and Ran(\(E\)) attack-heal}
\label{sec:compact-hall-ranE-attack-heal}

\providecommand{\Z}{\mathbb Z}
\providecommand{\C}{\mathbb C}
\providecommand{\Ran}{\operatorname{Ran}}
\providecommand{\Supp}{\operatorname{Supp}}
\providecommand{\Coh}{\operatorname{Coh}}
\providecommand{\Perf}{\operatorname{Perf}}
\providecommand{\Hilb}{\operatorname{Hilb}}
\providecommand{\Ext}{\operatorname{Ext}}
\providecommand{\RHom}{\operatorname{RHom}}
\providecommand{\BM}{\operatorname{BM}}
\providecommand{\Prim}{\operatorname{Prim}}
\providecommand{\sdim}{\operatorname{sdim}}
\providecommand{\smult}{\operatorname{smult}}
\providecommand{\Hall}{\mathcal H}
\providecommand{\Fact}{\mathsf{Fact}}
\providecommand{\Den}{\mathsf{Den}}
\providecommand{\Spch}{\mathrm{Sp}^{\mathrm{ch}}}
\providecommand{\hCS}{\mathrm{hCS}}
\providecommand{\van}{\varphi}

\subsection{Claim attacked}

The attacked claim is:
\[
\begin{gathered}
\hbox{the reduced OP/DT moduli problem on }X=S\times E\\
\Longrightarrow
\hbox{a compact Hall/factorization object on }\Ran(E)\\
\Longrightarrow
H^0\Prim_{\mathrm{prot}}\cong \mathfrak g_{\Delta_5}.
\end{gathered}
\]
This implication is false from the data presently in the manuscript.
The scalar quotient theory and the denominator algebra give the expected
Euler-characteristic shadow.  They do not construct the compact
Hall/factorization source.

The strongest honest statement available now is conditional.  If one
supplies an \(E\)-equivariant oriented critical Hall atlas before taking
the reduced \(E\)-quotient, a stability-sector HN completion, and a
projection-local or globally-wrapped Ran(\(E\)) locality theorem, then
one obtains an \(E\)-equivariant completed Hall factorization object on
\(\Ran(E)\).  If, in addition, protected integration gives
\[
\sdim\Prim_{\mathrm{prot},(a,k,b)}=f(ab,k)
\]
and the primitive Hall bracket satisfies the Gritsenko--Nikulin
Borcherds relations, then the resulting primitive Lie superalgebra has
denominator
\[
\operatorname{den}=64^{-1}\Delta_5(2Z).
\]
Those additional hypotheses are exactly the missing construction.

\subsection{Scalar moduli and the Hall moduli problem}

Let \(S\) be the elliptically fibered K3 surface used in the scalar
branch, let \(E\) be an elliptic curve, and set
\[
X=S\times E,\qquad p_E:X\longrightarrow E.
\]
For a primitive K3 class \(\beta_h\) with
\[
\langle\beta_h,\beta_h\rangle=2h-2
\]
and elliptic degree \(d\ge0\), the scalar coefficient spaces are
\[
\Hilb^\ell(X,(\beta_h,d))
=\{Z\subset X\mid [Z]=(\beta_h,d),\ \chi(\mathcal O_Z)=\ell\}
\]
and the stable-pair moduli spaces
\[
P_\ell(X,(\beta_h,d)).
\]
The elliptic curve acts by translation on the second factor:
\[
e\cdot(s,e')=(s,e+e').
\]
The reduced scalar invariants used in the manuscript are quotient
integrals
\[
\mathbf{DT}^{X}_{\ell,(\beta_h,d)}
=\int_{\Hilb^\ell(X,(\beta_h,d))/E}\nu\,de,
\qquad
\mathsf P^{X}_{\ell,(\beta_h,d)}
=\int_{P_\ell(X,(\beta_h,d))/E}\nu\,de,
\]
with \(\nu\) the Behrend function.  Oberdieck--Pixton,
Oberdieck--Pandharipande, and Oberdieck give the scalar chain recorded in
the manuscript:
\[
Z^X_{\mathrm{DT}}(T,Q,P)=Z^X_{\mathrm{OP}}(P,Q,T)
=-4096\,\Delta_5(Z)^{-2}.
\]

The Hall moduli problem is different.  It must be an extension-closed
derived stack of objects in a heart
\[
\mathcal A_\sigma\subset \Perf(X)
\]
with a stability condition \(\sigma\), HN property, support property, and
a charge map
\[
\operatorname{cl}:K_0(\mathcal A_\sigma)\longrightarrow
\Gamma_{\mathrm{ind}}=\Z^3.
\]
Write the Igusa charge as
\[
\gamma=(a,k,b),\qquad x^\gamma=q^a r^k s^b,
\]
so that the Borcherds degree is
\[
\alpha(a,k,b)=2a f_2-k f_3+2b f_{-2}
\]
and the expected signed primitive multiplicity is
\[
f(ab,k),\qquad
\phi_{0,1}(\tau,z)=\sum_{n,l}f(n,l)q^n r^l.
\]
The scalar variables identify \(a\) with the K3 norm exponent,
\(k\) with the D0/Fourier exponent, and \(b\) with the elliptic-degree
exponent after the OP chamber conversion.

\begin{proposition}[The scalar coefficient spaces are not the Hall stack]
\label{prop:scalar-spaces-not-ranE-hall}
The Hilbert-scheme and stable-pair quotient spaces used in the OP/DT
scalar theorem do not, by themselves, define an honest compact Hall
convolution object on \(\Ran(E)\).
\end{proposition}

\begin{proof}
Hall multiplication is defined by extension correspondences in an
extension-closed category:
\[
\mathfrak M_\gamma\times\mathfrak M_\eta
\xleftarrow{\ p\ }
\mathfrak E_{\gamma,\eta}
\xrightarrow{\ q\ }
\mathfrak M_{\gamma+\eta}.
\]
The class of ideal sheaves is not closed under arbitrary extensions:
an extension of two ideal sheaves need not be an ideal sheaf.  Likewise,
the stable-pair moduli space is a stable locus inside a larger tilted
heart; the stable-pair coefficient space is not itself the Hall stack on
which all extensions live.  Thus the scalar spaces are coefficient
spaces for a trace, not the extension-closed source of a Hall product.

The quotient by \(E\) creates a second obstruction.  The unreduced
extension stack is equivariant for simultaneous translation of the
subobject, quotient object, and middle object.  The product
\[
(\mathfrak M_\gamma/E)\times(\mathfrak M_\eta/E)
\]
has forgotten the relative \(E\)-position of the two inputs.  Extensions
depend on that relative position.  Therefore convolution on the separate
reduced quotients is not obtained by formal descent.  The Hall object
must be constructed equivariantly before reduced integration, or one must
construct an explicit reduced convolution kernel retaining the relative
position parameter.
\end{proof}

\subsection{The correct Ran(\(E\)) object}

The correct object is not a bare vector space and not the quotient scalar
series.  It is an \(E\)-equivariant, charge-completed, finite-first
factorization cosheaf or algebra on \(\Ran(E)\), built from an oriented
critical Hall stack before reduced integration.

For an open \(U\subset E\), the support-local candidate is
\[
\mathfrak M_{\sigma,S,\gamma}(U)
=
\{A\in\mathfrak M_{\sigma,S,\gamma}\mid
p_E(\Supp A)\subset U\}.
\]
At finite HN height \(R\), assuming critical charts and a strong
orientation \(o\), put
\[
\mathsf F_R\Hall^{\mathrm{or}}_{\sigma,S}(U)
=
\bigoplus_{\gamma\in\Gamma_{\sigma,S}^{HN}(R)}
H^{\BM}\!\left(
\mathfrak M_{\sigma,S,\gamma}(U),
\van_\gamma\otimes\mathscr L_{o,\gamma}
\right).
\]
The completed object is
\[
\widehat{\Hall}^{\mathrm{or}}_{\sigma,S}(U)
=
\varprojlim_R \mathsf F_R\Hall^{\mathrm{or}}_{\sigma,S}(U).
\]
Translation gives isomorphisms
\[
t_e^*\widehat{\Hall}^{\mathrm{or}}_{\sigma,S}(U)
\simeq
\widehat{\Hall}^{\mathrm{or}}_{\sigma,S}(U+e),
\]
so the object is naturally \(E\)-equivariant.  Reduced integration is a
final operation on this equivariant object.  It is not part of the
definition of the local factorization products.

\begin{theorem}[Conditional Ran(\(E\)) Hall construction]
\label{thm:conditional-ranE-hall-construction}
Assume the following data.
\begin{enumerate}
\item An extension-closed heart \(\mathcal A_\sigma\subset\Perf(X)\) with
a stability condition \(\sigma\), HN property, support property, and
bounded semistable substacks in strict sectors.
\item A charge map to \(\Gamma_{\mathrm{ind}}=\Z^3\) compatible with
\(\alpha(a,k,b)=2af_2-kf_3+2bf_{-2}\).
\item For every \(\gamma\), a derived critical stack
\(\mathfrak M_{\sigma,\gamma}\) with vanishing-cycle coefficients.
\item An \(E\)-equivariant strong orientation \(o\), compatible with
exact triangles and Thom--Sebastiani transport.
\item Extension correspondences
\[
\mathfrak M_\gamma\times\mathfrak M_\eta
\xleftarrow{\ p\ }
\mathfrak E_{\gamma,\eta}
\xrightarrow{\ q\ }
\mathfrak M_{\gamma+\eta}
\]
whose pull-push operations are admissible for the chosen
Borel--Moore/vanishing-cycle theory.
\item A projection-locality theorem: for disjoint opens
\(U,V\subset E\), objects supported over \(U\) have no extensions with
objects supported over \(V\), except through the direct-sum stratum.
\item A collision atlas over the diagonals in \(\Ran(E)\), identifying
the extension stacks in \(p_E^{-1}(D)\simeq S\times D\) with the chiral
collision operations.
\end{enumerate}
Then the assignment
\[
U\longmapsto \widehat{\Hall}^{\mathrm{or}}_{\sigma,S}(U)
\]
is an \(E\)-equivariant \(h_S\)-adically completed Hall factorization
object on \(\Ran(E)\).
\end{theorem}

\begin{proof}
At finite height \(R\), the support property and HN boundedness give a
finite set of surviving charges.  The \(h_S\)-height \(>R\) part is a
two-sided Hall ideal because \(h_S\) is positive and additive in the
strict sector.  Therefore \(\mathsf F_R\Hall^{\mathrm{or}}\) is a finite
Hall quotient and the transition maps form an inverse system.

For disjoint opens \(U,V\), projection-locality gives
\[
\Ext_X^*(A,B)=0
\]
whenever \(p_E(\Supp A)\subset U\) and \(p_E(\Supp B)\subset V\), apart
from the split direct-sum operation.  Thus the external product over
disjoint opens is the factorization product.  When points collide, the
extension stack over the diagonal supplies the chiral operation; the
strong orientation transports vanishing cycles by Thom--Sebastiani.
Associativity is the usual comparison over the stack of two-step
filtrations.  Passing to the inverse limit gives the completed object.
The \(E\)-equivariance follows from translation of supports and of the
extension correspondences.
\end{proof}

\subsection{Global-wrapping obstruction}

The preceding theorem is support-local.  The full OP scalar class is not
automatically support-local on \(E\).

\begin{lemma}[Positive elliptic degree is not Ran-support-local]
\label{lem:positive-e-degree-not-ran-local}
Let \(C\subset S\times E\) be a one-dimensional subscheme whose curve
class has positive \(E\)-degree.  Then \(p_E(C)\) is one-dimensional,
hence equals \(E\).  In particular \(C\) cannot be assigned to a proper
open \(U\subsetneq E\) by the condition \(p_E(C)\subset U\).
\end{lemma}

\begin{proof}
The pushforward of the fundamental one-cycle of \(C\) to \(E\) has
positive degree.  Hence the image of \(C\) under \(p_E\) is a closed
one-dimensional subset of the irreducible curve \(E\).  Therefore
\(p_E(C)=E\).
\end{proof}

\begin{corollary}[Naive support factorization misses the \(s\)-direction]
\label{cor:ranE-misses-s-direction}
The naive support-local Hall factorization object on \(\Ran(E)\) can
only see sectors whose support over \(E\) is contained in finite
subsets.  It does not directly realize the full Igusa product over
\[
(a,k,b)\in\Gamma_{\mathrm{eff}}
\]
when \(b\) is the OP elliptic-degree exponent and \(b>0\).
\end{corollary}

\begin{proof}
The \(s^b\)-direction records the elliptic-degree exponent in the OP
normalization.  By Lemma~\ref{lem:positive-e-degree-not-ran-local}, a
curve with positive elliptic degree dominates \(E\), so it is global from
the viewpoint of \(\Ran(E)\).  Disjoint-open factorization on
\(\Ran(E)\) cannot be defined by support separation for such objects.
\end{proof}

This is the sharp new obstruction.  A compact realization must choose one
of two routes.
\begin{enumerate}
\item Restrict the Hall factorization source to the projection-finite
sector and prove that the missing \(s\)-direction is produced by a
separate Fock/Hecke/global mode construction compatible with the
Borcherds product.
\item Enlarge the Ran(\(E\)) object to a hybrid factorization object:
local vertical insertions on \(\Ran(E)\), together with global
\(E\)-wrapping coefficient sectors.  The collision theory must then prove
that local collisions commute correctly with the global wrapped sector.
\end{enumerate}
Neither route is constructed in the present manuscript or the local
agent material.

\subsection{Collision locality}

For the projection-finite sector, locality is clean.  If
\(U_1,\ldots,U_r\subset E\) are pairwise disjoint, the factorization map
is induced by direct sum:
\[
\bigotimes_i
\widehat{\Hall}^{\mathrm{or}}_{\sigma,S}(U_i)
\longrightarrow
\widehat{\Hall}^{\mathrm{or}}_{\sigma,S}\!\left(\bigcup_iU_i\right).
\]
The reason is first-principles support separation: coherent objects with
disjoint supports on the smooth threefold have no nonzero local
\(\mathcal Ext\)-sheaves between them, hence no non-split extensions.

When points collide, the direct-sum stratum is replaced by the stack of
extensions
\[
\mathfrak E_{\gamma_1,\ldots,\gamma_r}(D)
\]
inside the formal tube \(p_E^{-1}(D)\simeq S\times D\), where \(D\) is a
small disc around the collision point.  The collision operation is the
pull-push with vanishing cycles and orientation transport:
\[
p^*(-)\xrightarrow{\mathrm{TS}_o}
\van_{\mathfrak E}\otimes\mathscr L_{o,\mathfrak E}
\xrightarrow{q_!}
\van_{\gamma_1+\cdots+\gamma_r}\otimes
\mathscr L_{o,\gamma_1+\cdots+\gamma_r}.
\]
Associativity is the comparison between the two ways of projecting from
the stack of two-step flags.  This is exactly the Hall source of the
chiral collision maps.  It is conditional on the critical atlas,
orientation, and admissible pull-push theory.

For globally \(E\)-wrapping sectors the locality statement is absent.
Such objects meet every nonempty open in \(E\).  They cannot be separated
by disjoint opens, and their extensions with local vertical objects are
not controlled by the usual Ran diagonal alone.  A compact realization
must supply an extension stack with one global leg and several local
legs, plus a proof that this hybrid operation satisfies the factorization
associativity constraints.

\subsection{Protected integration and denominator shadow}

Assume the conditional Ran(\(E\)) Hall object has been constructed.  A
protected integration is not just the Behrend-weighted scalar integral.
It must be a continuous character on the completed Hall object, compatible
with products, orientation, and the reduced \(E\)-integration:
\[
\chi^{\mathrm{prot}}:
K_0\Prim_{\mathrm{prot}}
\widehat{\Hall}^{\mathrm{or}}_{\sigma,S}
\longrightarrow
\widehat{\mathbb T}_{\Gamma_{\mathrm{eff}}}.
\]
The Igusa determinant follows only after the primitive equality
\[
\chi^{\mathrm{prot}}(\gamma)
=
\sdim\Prim_{\mathrm{prot},\gamma}
=f(ab,k),
\qquad \gamma=(a,k,b),
\]
and the orientation character comparison
\[
\epsilon_o=\nu_{\Delta_5}
\]
have been proved in the same normalization as the Borcherds product.
Then
\[
64(qrs)^{1/2}
\prod_{(a,k,b)\in\Gamma_{\mathrm{eff}}}
(1-q^a r^k s^b)^{f(ab,k)}
=\Delta_5(Z),
\]
and, in denominator variables,
\[
e^{-2\pi i(\rho,z)}
\prod_{\gamma\in\Gamma_{\mathrm{eff}}}
\left(1-e^{-2\pi i(\alpha(\gamma),z)}\right)^{f(ab,k)}
=64^{-1}\Delta_5(2Z).
\]
This proves only the determinant shadow.  The bracket theorem remains
separate:
\[
H^0\Prim_{\mathrm{prot}}
\widehat{\Hall}^{\mathrm{or}}_{\sigma,S}
\stackrel{?}{\cong}
\mathfrak g_{\Delta_5}.
\]
It requires the real simple primitives, Serre relations, imaginary
simple-root parities and multiplicities, imaginary orthogonality,
continuous Hopf pairing, and radical quotient.

\subsection{Failure theorem}

\begin{theorem}[Present obstruction]
\label{thm:ranE-present-obstruction}
The current manuscript and local agent material do not construct the
compact Hall/factorization object on \(\Ran(E)\) whose protected
primitive bracket is \(\mathfrak g_{\Delta_5}\).
\end{theorem}

\begin{proof}
The obstruction has five independent parts.

First, the OP/DT scalar branch is a quotient-by-\(E\) trace.  It is not
an extension-closed Hall stack.

Second, Hall convolution is equivariant before reduction.  Taking the
quotient by \(E\) separately on each input forgets relative position, so
the scalar quotient does not inherit convolution by formal descent.

Third, naive support-local factorization on \(\Ran(E)\) sees only
objects whose support projects to finite subsets of \(E\).  Positive
elliptic degree is global over \(E\), by
Lemma~\ref{lem:positive-e-degree-not-ran-local}.  Thus the full
\(s\)-direction of the Igusa product requires either a global-wrapping
enhancement or a separate Hecke/Fock construction.

Fourth, reduced obstruction theory and Behrend weighting compute scalar
Euler characteristics.  They do not by themselves supply a strong
orientation compatible with Hall extensions, reduced integration, and the
automorphic character \(\nu_{\Delta_5}\).

Fifth, the denominator identity determines signed root
supermultiplicities.  It does not determine the Hall primitive bracket or
the Borcherds defining relations.
\end{proof}

\subsection{Exact missing lemmas}

The compact construction closes only after the following statements are
proved.
\begin{enumerate}
\item \textbf{Heart and stability.}  Construct an extension-closed heart
\(\mathcal A_\sigma\subset\Perf(S\times E)\) with HN property, support
property, bounded semistable stacks, and charge map
\(\operatorname{cl}:K_0(\mathcal A_\sigma)\to\Z^3\) matching the Igusa
variables \((a,k,b)\).
\item \textbf{Projection-locality or global wrapping.}  Either prove that
the relevant compact BPS sector is projection-finite over \(E\), or build
the hybrid Ran(\(E\)) factorization object with global \(E\)-wrapping
coefficient sectors.
\item \textbf{Equivariant-before-reduction Hall descent.}  Construct the
\(E\)-equivariant Hall object before quotienting and the reduced
integration map after convolution.  Equivalently, construct the reduced
relative-position kernel.
\item \textbf{Critical atlas.}  Give local critical charts for the
derived moduli stacks and prove descent over the Cech/Ran nerve.
\item \textbf{Strong orientation.}  Construct an \(E\)-equivariant strong
orientation compatible with exact triangles and Thom--Sebastiani maps,
then prove \(\epsilon_o=\nu_{\Delta_5}\).
\item \textbf{HN finite-first completion.}  Prove the finite HN quotient
system and continuity of Hall multiplication in the \(h_S\)-adic topology.
\item \textbf{Reduced obstruction compatibility.}  Prove that the
reduced obstruction theory and Behrend/protected integration define a
Hall character, not merely scalar coefficients.
\item \textbf{Primitive multiplicities.}  Prove from the compact object
that
\[
\sdim\Prim_{\mathrm{prot},(a,k,b)}=f(ab,k),
\]
not by substituting the desired denominator.
\item \textbf{Borcherds bracket.}  Identify the real simple primitives
and prove the Serre relations, imaginary simple-root relations,
orthogonality relations, Hopf pairing, and radical quotient which give
\(\mathfrak g_{\Delta_5}\).
\item \textbf{Chiral comparison.}  Construct the conilpotent chiral bar
coalgebra \(C_X\simeq\bar B_E^{\mathrm{ch}}(\mathcal F_X)\) and a
quasi-isomorphism
\[
\Omega_E^{\mathrm{ch}}C_X\simeq
\Den_{\Delta_5,E}.
\]
\end{enumerate}

\subsection{Claim-status recommendation}

Keep the compact \(E_3\)/Hall/chiral realization as an open problem.
Strengthen the formulation by adding two precise obstructions:
\[
\hbox{Hall before }E\hbox{-reduction},
\qquad
\hbox{projection-local Ran(\(E\)) versus global \(E\)-wrapping sectors}.
\]
The conditional theorem
\ref{thm:conditional-ranE-hall-construction} is safe to use as the
maximal construction statement.  It is a theorem only from the named
atlas, orientation, locality, and completion hypotheses.  The full
compact Igusa realization is not proved.

\subsection{Anchors}

Local manuscript anchors:
\begin{itemize}
\item \verb|main.tex:55--117|: abstract status; determinant and
denominator proved, compact Hall/chiral source isolated as realization
problem.
\item \verb|main.tex:250--302|: expected \(E_3\)-source,
\(\Spch_{K3,E}\), \(\mathcal F_X\in\Fact(\Ran(E))\), and primitive
comparison with \(\mathfrak g_{\Delta_5}\).
\item \verb|main.tex:721--743|: scalar Hilbert-scheme quotient by \(E\)
and Behrend weighting.
\item \verb|main.tex:840--887|: reduced Hilbert-scheme DT/PT/OP scalar
square and constant \(-4096\).
\item \verb|main.tex:1673--1716|: formal
\(\Den_{\Delta_5,E}=U_E^{\mathrm{ch}}(\mathfrak g_{\Delta_5}\otimes
\mathcal O_E)\).
\item \verb|main.tex:1718--1808|: \(E_3\)-to-chiral specialization and
Koszul comparison are required data, not consequences.
\item \verb|main.tex:1834--1905|: sectorial HN truncation and finite
Hall quotient.
\item \verb|main.tex:1907--2096|: compact Igusa realization datum,
criterion, determinant-does-not-determine-Hall-constants, and open
problem.
\item \verb|main.tex:2436--2516|: main theorem; determinant and
denominator are proved, compact microscopic comparison remains open.
\item \verb|agent_material/10_compact_hall_atlas_attack_heal.tex:220|:
the reduced \(E\)-quotient does not automatically carry Hall convolution.
\item \verb|agent_material/11_compact_hall_finite_heart_and_hn.tex:293|:
the correct finite-first object is the HN quotient system.
\item \verb|agent_material/11_factorization_koszul_minimal_data.tex:427|:
the five compact source/comparison blocks are not presently built.
\end{itemize}

Primary source anchors already present in \verb|proj.bib|:
\begin{itemize}
\item \cite{OB}, Theorem 1 and product (3): OP scalar branch and Igusa
product normalization.
\item \cite{OPand}: reduced GW/PT comparison for primitive K3 classes on
\(S\times E\).
\item \cite{OberdieckReducedPT}, Theorem 3(i): reduced DT/PT comparison.
\item \cite{BRYAN}, Definition 2.1 and Conjecture 2.3: quotient
Hilbert-scheme DT normalization and Bryan variable convention.
\item \cite{DT}, \cite{BEHREND}: DT obstruction theory and Behrend
function.
\item \cite{KSCoHA}, \cite{DavisonMeinhardtCoHA}: critical/cohomological
Hall technology once the moduli, orientation, and correspondences are
supplied.
\item \cite{BD}: chiral envelope and chiral collision formalism on a
curve.
\item \cite{CostelloRenormalization}, \cite{CostelloGwilliam1},
\cite{CostelloGwilliam2}: BV observables give factorization algebras
after QME, locality, descent, and anomaly hypotheses.
\item \cite{B}, \cite{GN}: Borcherds product and
Gritsenko--Nikulin denominator algebra.
\end{itemize}

\subsection{Files changed and checks}

Files changed:
\begin{center}
\texttt{agent\_material/12\_compact\_hall\_ranE\_attack\_heal.tex}.
\end{center}

Checks run in this attack-heal pass:
\begin{itemize}
\item Read \verb|CLAUDE.md|, \verb|AGENTS.md|,
\verb|/Users/raeez/ecosystem/INVARIANTS.md|, and the relevant
\verb|main.tex| loci before writing.
\item Grepped \verb|main.tex|, \verb|proj.bib|, and existing
\verb|agent_material| files for Hall, Ran, compact realization,
reduced quotient, Behrend, factorization, and denominator anchors.
\item Ran \verb|git diff --no-index --check /dev/null| on this file;
no whitespace errors were reported.
\item Ran \verb|git diff --cached --check| after staging this file; no
whitespace errors were reported.
\item No manuscript build was run: \verb|main.tex| was not edited, and
this file is not included in the TeX root.
\end{itemize}

\subsection{Remaining open questions}

The decisive open question is not whether the scalar square is correct;
it is correct in the OP/DT normalization recorded in the manuscript.  The
decisive question is whether one can construct a compact
\(E\)-equivariant Hall/factorization source whose reduced protected trace
is that scalar square and whose primitive bracket is
\(\mathfrak g_{\Delta_5}\).

The first obstruction to attack next is the global-wrapping problem:
positive elliptic degree is global over \(E\), while naive
Ran(\(E\))-factorization is support-local.  A proof must either remove
that sector from the local Hall source and recover it by a separate
Hecke/Fock mechanism, or construct the hybrid local/global collision
stack and prove its associativity and orientation compatibility.
